\documentclass[11pt]{article}
\usepackage{amsmath,amssymb,amsthm}
\usepackage[margin=1.1in]{geometry}
\usepackage{booktabs}

\newtheorem{proposition}{Proposition}
\newtheorem{remark}{Remark}

\title{Loss Function Derivatives:\\
Every Gradient the Repository Descends, Derived From Scratch}
\author{Yuval Lavie}
\date{\today}

\begin{document}
\maketitle

\noindent
The losses page (\texttt{Theory/losses}) catalogues what each loss
\emph{remembers}; this page computes what each loss \emph{does} during
training --- its gradient, derived step by step with the rules page
(\texttt{rules.tex}) as the toolkit and the activation-functions page
supplying $\sigma'$ and the softmax Jacobian. The punchline, reached
three separate times, is the same beautiful formula:
\emph{gradient = prediction $-$ target}.

\section{Squared error}

Per example, with prediction $\hat y$ and target $y$ (the $\tfrac12$ is
cosmetic, chosen to cancel the exponent):
\begin{equation}
  \ell(\hat y) = \tfrac{1}{2}\,(\hat y - y)^{2}.
  \label{eq:mse}
\end{equation}

\paragraph{Derivation.} Chain rule with outer $u \mapsto \tfrac12 u^2$
(power rule) and inner $u = \hat y - y$ (derivative $1$):
\begin{equation}
  \boxed{\;\frac{d\ell}{d\hat y} = \hat y - y\;}
  \label{eq:mseprime}
\end{equation}
Prediction minus target --- the archetype. For a linear model
$\hat y = w^{\top}x + b$, one more chain-rule step gives
$\nabla_w \ell = (\hat y - y)\,x$ and
$\partial \ell/\partial b = \hat y - y$: the residual, redistributed
along the input. Summed over data this is the normal-equations gradient
of the linear-regression note (rules page, eq.\ for
$\nabla \lVert y - Xw \rVert^2$).

\section{Absolute error and Huber}

\begin{equation}
  \ell_{\mathrm{abs}}(\hat y) = |\hat y - y|
  \qquad\Longrightarrow\qquad
  \frac{d\ell_{\mathrm{abs}}}{d\hat y}
  = \operatorname{sign}(\hat y - y),
  \quad
  \partial \ell_{\mathrm{abs}}\big|_{\hat y = y} = [-1, 1].
  \label{eq:maeprime}
\end{equation}
Piecewise: for $\hat y > y$ the loss is $\hat y - y$ with slope $+1$;
for $\hat y < y$ it is $y - \hat y$ with slope $-1$; at the kink the
subdifferential $[-1,1]$ applies (rules page). The gradient's
\emph{magnitude never exceeds 1} regardless of the residual --- that
bounded pull is the robustness of the median, in gradient form.

\paragraph{Huber.} With residual $r = \hat y - y$ and threshold
$\delta$:
\begin{equation}
  \ell_{\delta}(r) =
  \begin{cases}
    \tfrac12 r^{2} & |r| \le \delta\\[2pt]
    \delta\bigl(|r| - \tfrac{\delta}{2}\bigr) & |r| > \delta
  \end{cases}
  \qquad\Longrightarrow\qquad
  \frac{d\ell_{\delta}}{dr} =
  \begin{cases}
    r & |r| \le \delta\\[2pt]
    \delta\operatorname{sign}(r) & |r| > \delta.
  \end{cases}
  \label{eq:huberprime}
\end{equation}
Each branch differentiates by the rules above; at $|r| = \delta$ the
two branches agree ($r = \pm\delta$ vs.\
$\delta\operatorname{sign}(r) = \pm\delta$), so the derivative is
continuous --- Huber is $C^1$ by construction: it is squared error's
gradient, \emph{clipped} at $\pm\delta$. (Gradient clipping in deep
learning is exactly a Huberization applied after the fact.)

\section{Binary cross-entropy}

With label $y \in \{0,1\}$ and predicted probability $p \in (0,1)$:
\begin{equation}
  \ell(p) = -\,y\ln p - (1-y)\ln(1-p).
  \label{eq:bce}
\end{equation}

\paragraph{Derivative with respect to $p$.} Linearity plus the
log-derivative rule on each term
($\frac{d}{dp}\ln p = \frac1p$;
$\frac{d}{dp}\ln(1-p) = \frac{-1}{1-p}$ by the chain rule):
\begin{equation}
  \frac{d\ell}{dp}
  = -\frac{y}{p} + \frac{1-y}{1-p}
  = \frac{-y(1-p) + (1-y)p}{p(1-p)}
  = \frac{p - y}{p\,(1-p)}.
  \label{eq:bceprimep}
\end{equation}
Already ``prediction minus target,'' but divided by the Bernoulli
variance $p(1-p)$ --- which blows up as $p \to 0$ or $1$. Training on
$p$ directly would be numerically vicious. The cure is to differentiate
with respect to the \emph{logit} instead.

\paragraph{Derivative with respect to the logit $z$, where
$p = \sigma(z)$.} Chain rule, composing \eqref{eq:bceprimep} with
$\sigma'(z) = \sigma(z)(1 - \sigma(z)) = p(1-p)$ from the
activation-functions page:
\begin{align*}
  \frac{d\ell}{dz}
  = \frac{d\ell}{dp}\cdot\frac{dp}{dz}
  = \frac{p - y}{p(1-p)} \cdot p(1-p)
  = p - y.
\end{align*}
\begin{equation}
  \boxed{\;\frac{d\ell}{dz} = \sigma(z) - y\;}
  \label{eq:bceprimez}
\end{equation}
The Bernoulli variance cancels \emph{exactly}. This is not luck: the
sigmoid is the canonical link of the Bernoulli exponential family, and
for canonical links the NLL gradient in the natural parameter is always
mean minus observation. It is why logistic regression's gradient
(logistic-regression note) looks identical to linear regression's, and
why \texttt{BCEWithLogitsLoss} exists: computing
\eqref{eq:bceprimez} directly never forms the unstable quotient
\eqref{eq:bceprimep}.

\section{Categorical cross-entropy through softmax}

Labels one-hot $y \in \{0,1\}^{K}$ (with $\sum_k y_k = 1$), logits
$z \in \mathbb{R}^{K}$, probabilities $s = \operatorname{softmax}(z)$:
\begin{equation}
  \ell(z) = -\sum_{k=1}^{K} y_k \ln s_k(z) = -\ln s_c(z)
  \quad\text{for the true class } c.
  \label{eq:ce}
\end{equation}

\paragraph{Derivation.} Chain rule through the softmax Jacobian
$\partial s_i/\partial z_j = s_i(\delta_{ij} - s_j)$
(activation-functions page), with the log-derivative rule on each
$\ln s_k$:
\begin{align*}
  \frac{\partial \ell}{\partial z_j}
  = -\sum_{k} \frac{y_k}{s_k}\cdot\frac{\partial s_k}{\partial z_j}
  = -\sum_{k} \frac{y_k}{s_k}\; s_k\,(\delta_{kj} - s_j)
  = -\sum_{k} y_k\,(\delta_{kj} - s_j)
  = -\,y_j + s_j \sum_{k} y_k .
\end{align*}
Since $\sum_k y_k = 1$:
\begin{equation}
  \boxed{\;\nabla_z\,\ell = s - y\;}
  \label{eq:ceprime}
\end{equation}
Prediction minus target, for the third time --- softmax is the
canonical link of the categorical family, so the same
exponential-family cancellation fires. Note what \emph{didn't} happen:
we never materialized the $K \times K$ Jacobian; the $1/s_k$ from the
log and the $s_k$ from the Jacobian annihilate. This collapse is why
every framework fuses softmax and cross-entropy into one op
(\texttt{CrossEntropyLoss} in the softmax-regression and
\texttt{llm/} notes): the fused gradient \eqref{eq:ceprime} is stable
where the two-step computation is not. The derivation holds verbatim
for soft labels ($y$ any distribution), which is the label-smoothing
and distillation case.

\section{KL divergence}

For a fixed target distribution $q$ and model $s(z)$,
$D_{\mathrm{KL}}(q\,\Vert\,s) = \sum_k q_k \ln\frac{q_k}{s_k}
= \underbrace{\textstyle\sum_k q_k \ln q_k}_{\text{const in } z}
- \sum_k q_k \ln s_k$. The variable part is cross-entropy with soft
labels $q$, so by \eqref{eq:ceprime}:
\begin{equation}
  \nabla_z\, D_{\mathrm{KL}}\bigl(q \,\Vert\, s(z)\bigr) = s - q .
  \label{eq:klprime}
\end{equation}
Minimizing KL to a target and minimizing cross-entropy against it are
the same descent --- the information-theory note's identity
$H(q, s) = H(q) + D_{\mathrm{KL}}(q\Vert s)$, now in gradient form.

\section{Margin losses: hinge, perceptron, logistic, exponential}

Binary labels $y \in \{-1, +1\}$, score $f$, margin $m = y f$. Each
loss is $\ell(m)$; by the chain rule
$\frac{\partial \ell}{\partial f} = \ell'(m)\, y$, so it suffices to
differentiate in $m$.

\paragraph{Hinge.} $\ell(m) = \max(0, 1 - m)$. Piecewise: slope $-1$
where $m < 1$, slope $0$ where $m > 1$, subdifferential $[-1, 0]$ at
the kink $m = 1$:
\begin{equation}
  \frac{d\ell}{dm} =
  \begin{cases}
    -1 & m < 1\\
    0 & m > 1,
  \end{cases}
  \qquad\Longrightarrow\qquad
  \nabla_w \ell = -\,y\,x\;\mathbf{1}[\,y\,w^{\top}x < 1\,]
  \label{eq:hingeprime}
\end{equation}
for a linear scorer $f = w^{\top}x$. The gradient is \emph{all or
nothing}: examples with comfortable margin contribute exactly zero ---
this hard zero is where support-vector sparsity comes from (svm note).
The perceptron loss $\max(0, -m)$ differentiates identically with the
kink at $0$: update $y\,x$ on mistakes, nothing otherwise --- the
perceptron algorithm \emph{is} subgradient descent on this loss.

\paragraph{Logistic.} $\ell(m) = \ln(1 + e^{-m})$. Log-derivative rule
with $f(m) = 1 + e^{-m}$, then multiply top and bottom by $e^{m}$:
\begin{equation}
  \frac{d\ell}{dm}
  = \frac{-e^{-m}}{1 + e^{-m}}
  = -\,\frac{1}{e^{m} + 1}
  = -\,\sigma(-m).
  \label{eq:logisticprime}
\end{equation}
A \emph{soft} version of hinge's indicator: instead of switching off at
$m = 1$, the pull decays smoothly as $\sigma(-m) \to 0$ for
confidently-correct examples --- consistent with
\eqref{eq:bceprimez}, since this is the same loss in $\pm 1$
convention. Every example keeps a nonzero gradient; no sparsity, full
calibration.

\paragraph{Exponential.} $\ell(m) = e^{-m}$, so
$\ell'(m) = -e^{-m}$: the pull on a misclassified example grows
\emph{exponentially} with how wrong it is. That unbounded appetite for
hard examples is AdaBoost's reweighting rule --- and its notorious
outlier sensitivity --- in one line.

\section{Summary table}

Gradients with respect to the quantity actually optimized:

\begin{center}
\begin{tabular}{lll}
\toprule
Loss & Gradient & Shape of the pull\\
\midrule
Squared ($\tfrac12 r^2$) & $\hat y - y$ & linear in residual\\
Absolute & $\operatorname{sign}(\hat y - y)$ & constant magnitude\\
Huber & $r$ clipped to $[-\delta, \delta]$ & linear, then capped\\
BCE on logit & $\sigma(z) - y$ & bounded by $1$\\
Softmax CE on logits & $s - y$ & bounded, per class\\
KL to target $q$ & $s - q$ & same as CE\\
Hinge & $-y\,\mathbf{1}[m < 1]$ & on/off\\
Perceptron & $-y\,\mathbf{1}[m < 0]$ & on/off at $0$\\
Logistic (margin form) & $-y\,\sigma(-m)$ & smooth on/off\\
Exponential & $-y\,e^{-m}$ & unbounded\\
\bottomrule
\end{tabular}
\end{center}

\medskip
\noindent
Three families, one moral. The probabilistic losses (squared, BCE,
softmax CE) all reduce to \emph{prediction minus target} in the right
parameterization --- the exponential-family/canonical-link
cancellation, witnessed three times above. The robust losses (absolute,
Huber) are the same pull with its magnitude capped. The margin losses
differ only in how fast the pull switches off as the margin grows:
never (exponential), smoothly (logistic), abruptly at $1$ (hinge), or
at $0$ (perceptron) --- and that switch-off schedule is precisely what
decides calibration versus sparsity on the losses page.

\end{document}
